\section{Introduction to Artificial Intelligence}
\label{sec:ai-introduction}
\slidesrc{Thinking machines: 1--97}

Before studying logic, Prolog and agents, the course opens with a big-picture tour: what
\emph{intelligence} is, how \emph{artificial} intelligence (AI) was born and grew, what it can
do today, why people are worried about it, and where this course sits in the AI landscape.
This chapter covers:
\begin{enumerate}[nosep]
  \item what intelligence is (several definitions, Kahneman's System~1 and System~2);
  \item what artificial intelligence is, and why it is defined through \emph{agents};
  \item a short history of AI, from Aristotle to self-driving cars;
  \item recent achievements (AlphaStar, AlphaFold, generative and agentic AI);
  \item the risks, and the principles and rules written to control them;
  \item symbolic vs.\ sub-symbolic AI, and why this course is about the \emph{symbolic}
        and \emph{distributed} side.
\end{enumerate}

\subsection{About the course}
\label{sec:about-the-course}
\slidesrc{Thinking machines: 1--3}

The course \textbf{Symbolic and Distributed Artificial Intelligence} (SDAI, sometimes
written S\&DAI) is taught by prof.\ Viviana Mascardi at the University of Genoa. All
material (slides, announcements, lab instructions) is published on the course page of the
university e-learning platform (AulaWeb).

The course has two souls, already visible in its name:
\begin{itemize}[nosep]
  \item \textbf{Symbolic}: intelligence obtained by manipulating \emph{symbols} with explicit
        rules: logic (Part~II of these notes) and the logic programming language
        \textbf{Prolog} (Part~IV);
  \item \textbf{Distributed}: intelligence spread over many autonomous software
        \textbf{agents} that interact: agents, multi-agent systems, their architectures,
        their communication and the platforms used to build them (Parts~III and~V).
\end{itemize}
Section~\ref{sec:symbolic-vs-subsymbolic} explains, at the end of this chapter, why these
two things belong together.

\subsection{What is intelligence? What is AI?}
\label{sec:what-is-intelligence}
\slidesrc{Thinking machines: 4--12}

To build \emph{artificial} intelligence we should first agree on what intelligence is.
There is no single accepted answer; the slides show several complementary views.

\subsubsection{The word and a classic definition}
\slidesrc{Thinking machines: 4--6}

The word comes from Latin \emph{intellig\u{e}re}: ``to understand'', literally
``to read inside'' (\emph{intus legere}) or ``to read between'' (\emph{inter legere}),
that is, to read beyond the surface of things.

\begin{nbdef}{Intelligence (Mainstream Science on Intelligence, 1994)}
Intelligence is a very general mental capability that, among other things, involves the
ability to \textbf{reason}, \textbf{plan}, \textbf{solve problems}, \textbf{think abstractly},
\textbf{comprehend complex ideas}, \textbf{learn quickly} and \textbf{learn from experience}.
It is not merely book-learning, a narrow academic skill, or test-taking smarts. Rather, it
reflects a broader and deeper capability for comprehending our surroundings:
``catching on'', ``making sense'' of things, or ``figuring out'' what to do.
\end{nbdef}

The definition was published in the \emph{Wall Street Journal} in December 1994 and signed
by 52 scientists. The professor remarks that the statement also contains claims she does not
feel comfortable with: it is a reference point, not an undisputed truth.
The seven abilities in bold (highlighted in red on slide~6) are worth remembering: each of
them has become a branch of AI.

\begin{center}
\small
\begin{tabular}{@{}ll@{}}
\toprule
\textbf{Ability in the definition} & \textbf{AI field that tries to reproduce it} \\
\midrule
reason                          & logic, automated reasoning, Prolog (this course) \\
plan                            & automated planning (e.g.\ the STRIPS planner of Shakey) \\
solve problems                  & search algorithms, constraint solving \\
think abstractly                & knowledge representation, ontologies \\
comprehend complex ideas        & natural language understanding \\
learn quickly / from experience & machine learning, reinforcement learning \\
\bottomrule
\end{tabular}
\captionof{table}{The abilities in the 1994 definition and the AI research areas that study them.}
\end{center}

\subsubsection{Many kinds of intelligence}
\slidesrc{Thinking machines: 7--8}

Intelligence is not a single number. Two influential theories:

\begin{nbdef}{Multiple intelligences (H.~Gardner, \emph{Frames of Mind}, 1983)}
Gardner proposed that humans have several distinct intelligences:
\begin{multicols}{2}
\begin{enumerate}[nosep]
  \item musical-rhythmic and harmonic;
  \item visual-spatial;
  \item verbal-linguistic;
  \item logical-mathematical;
  \item bodily-kinesthetic;
  \item interpersonal (understanding others);
  \item intrapersonal (understanding oneself);
  \item naturalistic (recognising and classifying nature);
  \item existential (asking deep questions about life).
\end{enumerate}
\end{multicols}
\end{nbdef}

\begin{nbdef}{Emotional intelligence (D.~Goleman, 1995)}
Emotional intelligence (EI) is the capability of individuals to recognise their own emotions
and those of others, discern between different feelings and label them appropriately, use
emotional information to guide thinking and behaviour, and manage and/or adjust emotions
to adapt to environments or achieve one's goal(s).
\end{nbdef}

\begin{nbintuition}[Why does an AI course care?]
If intelligence has many faces, then ``artificial intelligence'' can also mean many
different things. A chess program shows logical-mathematical intelligence; a robot that
walks on rough ground shows bodily-kinesthetic intelligence; a chatbot that adapts its tone
to an angry customer imitates a bit of emotional and interpersonal intelligence. A
\emph{multi-agent system} (the main topic of this course) needs \emph{interpersonal}
(social) intelligence: agents must understand and negotiate with each other.
\end{nbintuition}

\subsubsection{Thinking, fast and slow: System 1 and System 2}
\label{sec:system1-system2}
\slidesrc{Thinking machines: 9--11}

In the book \emph{Thinking, Fast and Slow} (2011), the psychologist and Nobel prize winner
Daniel Kahneman describes the mind as two cooperating ``systems''.

\begin{nbdef}{System 1 and System 2}
\begin{itemize}[nosep]
  \item \textbf{System 1}: fast, automatic, frequent, emotional, stereotypic, unconscious.
  \item \textbf{System 2}: slow, effortful, infrequent, logical, calculating, conscious.
\end{itemize}
\end{nbdef}

\begin{center}
\small
\begin{tabularx}{\linewidth}{@{}XX@{}}
\toprule
\textbf{System 1 can\ldots{}} (in order of complexity) & \textbf{System 2 is needed to\ldots{}} \\
\midrule
determine that an object is farther than another & prepare for the start of a sprint \\
localise the source of a sound & direct your attention to the clowns at the circus \\
complete the phrase ``war and \ldots'' & look for the woman with grey hair \\
show disgust at a gruesome image & try to recognise a sound \\
solve $2+2=?$ & sustain a faster-than-normal walking rate \\
read text on a billboard & judge if a behaviour is appropriate in a social setting \\
drive a car on an empty road & count the number of A's in a text \\
think of a good chess move (if you are a chess master) & give someone your phone number \\
understand simple sentences & park into a tight parking space \\
 & compare the price/quality ratio of two washing machines \\
 & determine the validity of a complex logical reasoning \\
\bottomrule
\end{tabularx}
\captionof{table}{Kahneman's examples (slides 10--11). Notice that the same task can move
from System~2 to System~1 with practice: a chess master ``sees'' a good move.}
\end{center}

\begin{nbexample}{the bat and the ball}
\notinslides
A bat and a ball cost 1.10 euro in total. The bat costs 1 euro more than the ball. How much
does the ball cost? System~1 shouts ``10 cents!''. System~2, if you wake it up, checks:
if the ball costs $x$, the bat costs $x+1$, so
\begin{align*}
x + (x + 1) &= 1.10 \\
2x &= 0.10 \\
x &= 0.05 .
\end{align*}
The ball costs 5 cents (check: $0.05 + 1.05 = 1.10$). Fast intuition is cheap but can be
wrong; slow reasoning is expensive but verifiable. Keep this in mind: it is exactly the
difference between \emph{sub-symbolic} and \emph{symbolic} AI
(Section~\ref{sec:symbolic-vs-subsymbolic}).
\end{nbexample}

\subsubsection{Artificial intelligence and intelligent agents}
\label{sec:ai-definition}
\slidesrc{Thinking machines: 12}

\begin{nbdef}{Artificial intelligence}
Artificial intelligence (AI) is intelligence demonstrated by machines, unlike the natural
intelligence displayed by humans and animals. The leading AI textbook, \emph{Artificial
Intelligence: A Modern Approach} (S.~Russell and P.~Norvig, 4th edition), defines the field
as the study of \textbf{intelligent agents}: any device that \textbf{perceives} its
environment and \textbf{takes actions} that maximise its chance of successfully achieving
its \textbf{goals}.
\end{nbdef}

This definition is the bridge to the whole course: AI is not defined by \emph{how} the
device works inside (neural network, logic rules, \ldots) but by its behaviour as an
\textbf{agent}, a loop of perceiving and acting towards goals.

\begin{center}
\begin{tikzpicture}[>=Stealth, node distance=1.2cm,
  box/.style={draw, rounded corners, thick, minimum width=2.6cm, minimum height=1cm, align=center}]
  \node[box, draw=nbblue, fill=nbblue!8] (agent) {\textbf{Agent}\\[-1pt]\footnotesize goals};
  \node[box, draw=nbgreen, fill=nbgreen!8, right=3.2cm of agent] (env) {\textbf{Environment}};
  \draw[->, thick, nbgreen!70!black] (env.north west) to[bend right=25] node[above, font=\small]{percepts (sensors)} (agent.north east);
  \draw[->, thick, nbred] (agent.south east) to[bend right=25] node[below, font=\small]{actions (actuators)} (env.south west);
\end{tikzpicture}
\captionof{figure}{An intelligent agent: it perceives the environment and acts on it to reach its
goals. Part~III of these notes is entirely about this picture.}
\end{center}

\begin{nbwhy}
Everyday intelligent agents: a robot vacuum cleaner (perceives obstacles and dirt, acts by
moving and sucking, goal: a clean floor); a smart thermostat (perceives temperature and your
habits, acts on the boiler, goal: comfort at minimum cost); a trading bot (perceives prices,
acts by buying/selling, goal: profit); a navigation app (perceives traffic, acts by
suggesting a route, goal: get you there fast).
\end{nbwhy}

\subsection{A short history of AI}
\label{sec:ai-history}
\slidesrc{Thinking machines: 13--67}

The slides build a timeline one event at a time (slides 16--67). The complete timeline is
redrawn in Figure~\ref{fig:ai-timeline}; the sub-sections below explain each step.

\begin{center}
\begin{tikzpicture}[font=\small,
  ev/.style={anchor=north west, text width=10.6cm, align=left, inner sep=2pt},
  yr/.style={anchor=north east, font=\small\bfseries, inner sep=2pt}]
  \coordinate (e0) at (0.3,0);
  \foreach \i/\year/\txt/\c in {
    1/1950/{Alan Turing proposes the ``Turing test'' to determine a machine's human-like ability}/nbblue,
    2/1956/{The term ``Artificial Intelligence'' is coined at the Dartmouth workshop}/nbblue,
    3/1966--72/{Stanford Research Institute builds Shakey, the first robot that reasons about its own actions}/nbgreen,
    4/1970--72/{Colmerauer designs Prolog; first Prolog system by Colmerauer and Roussel}/nbgreen,
    5/1974--80/{First ``AI winter'': the Lighthill report (1973) cuts funding}/nbred,
    6/1980/{Symbolics Lisp machines are commercialised: AI renaissance (expert systems)}/nbgreen,
    7/1987--93/{Second ``AI winter'': collapse of the Lisp machine market}/nbred,
    8/1997/{IBM Deep Blue defeats world chess champion Garry Kasparov}/nbgreen,
    9/2008/{Google speech-to-search iPhone app}/nborange,
    10/2011/{IBM Watson defeats Ken Jennings at the quiz show Jeopardy!}/nborange,
    11/2011--12/{Apple introduces Siri, an intelligent personal assistant}/nborange,
    12/2012/{Google Brain learns to identify cats on YouTube (75\% accuracy) without being told what a cat is}/nborange,
    13/2013/{China's Tianhe-2 doubles the world's top supercomputing speed (33.86 petaflops)}/nbgray,
    14/2014/{Facebook develops DeepFace, face recognition with near-human accuracy}/nborange,
    15/2015/{Google open-sources the TensorFlow machine learning library}/nborange,
    16/2016/{Google DeepMind's AlphaGo defeats Lee Sedol at Go}/nborange,
    17/2016/{China's Sunway TaihuLight triples the top supercomputing speed (93 petaflops)}/nbgray,
    18/2016/{Uber pilots a self-driving car programme in Pittsburgh, PA}/nborange,
    19/2019--now/{AlphaStar, AlphaFold, generative AI, agentic AI (Section~\ref{sec:recent-achievements})}/nbpurple}
  {
    \pgfmathtruncatemacro{\prev}{\i-1}
    \node[ev] (e\i) at ([yshift=-3pt]e\prev.south west) {\txt};
    \node[yr, text=\c] at ([xshift=-0.55cm]e\i.north west) {\year};
    \fill[\c] ([xshift=-0.3cm, yshift=-7pt]e\i.north west) circle (3.2pt);
  }
  \begin{scope}[on background layer]
    \draw[very thick, nbblue!50] (0,0) -- (0,0 |- e19.south);
  \end{scope}
\end{tikzpicture}
\captionof{figure}{The AI timeline of the slides (slide 67, extended). Colours:
\textcolor{nbblue}{foundations}, \textcolor{nbgreen}{symbolic AI successes},
\textcolor{nbred}{AI winters}, \textcolor{nborange}{data-driven / neural successes},
\textcolor{nbgray}{computing power}, \textcolor{nbpurple}{today}. Notice the shift: until the
1990s the big results are symbolic (reasoning, planning, chess search); after 2008 they are
mostly learned from data.}
\label{fig:ai-timeline}
\end{center}

\subsubsection{Before computers}
\slidesrc{Thinking machines: 13--15}

The slides show three portraits, without text. Here is why each person is on the
``before'' list (the explanations are added):
\begin{itemize}[nosep]
  \item \textbf{Aristotle} (384--322 BC) invented the first system of formal logic, the
        \emph{syllogism}: ``All men are mortal; Socrates is a man; therefore Socrates is
        mortal''. The very same example will come back in first-order logic
        (Whiteboard~\ref{wb:modus-ponens}) and in Prolog.
  \item \textbf{Gottfried Wilhelm von Leibniz} (1646--1716) dreamed of a universal symbolic
        language (\emph{characteristica universalis}) and a ``calculus of reasoning''
        (\emph{calculus ratiocinator}), so that disputes could be settled by saying
        ``let us calculate!''. He also studied binary arithmetic and built a mechanical
        calculator. This is the dream of symbolic AI: reasoning as computation.
  \item \textbf{Mary Shelley} (1797--1851) wrote \emph{Frankenstein} (1818), the story of an
        artificial being created by a scientist, and of the consequences of creating it:
        the first great story about the risks of artificial creatures
        (see Section~\ref{sec:risks-and-regulation}).
\end{itemize}

\subsubsection{1950: the Turing test}
\slidesrc{Thinking machines: 16--23}

In 1950 Alan Turing published the paper \emph{Computing Machinery and Intelligence} in the
philosophy journal \emph{Mind}. Instead of answering the vague question ``Can machines
think?'', he replaced it with a game that can actually be played, the
\textbf{imitation game}, now called the \textbf{Turing test}.

\begin{nbdef}{Turing test}
A human interrogator chats (in writing) with two hidden partners: a human and an AI system.
By asking questions, the interrogator must decide which one is the machine. If the
interrogator cannot reliably tell them apart, the machine passes the test.
\end{nbdef}

\begin{center}
\begin{tikzpicture}[>=Stealth, font=\small,
  who/.style={draw, rounded corners, thick, minimum width=2.4cm, minimum height=0.8cm, align=center}]
  \node[who, draw=nbblue, fill=nbblue!8] (int) at (0,0) {Human\\interrogator};
  \draw[ultra thick, nbgray] (2.2,-1.6) -- (2.2,1.6) node[above, font=\scriptsize]{wall};
  \node[who, draw=nbgreen, fill=nbgreen!8] (h) at (5,0.85) {Human (A or B?)};
  \node[who, draw=nbred, fill=nbred!8] (m) at (5,-0.85) {AI system (A or B?)};
  \draw[<->, thick] (int.east) -- (2.2,0) node[right, font=\Large]{?};
  \draw[<->, thick] (2.5,0) -- (h.west);
  \draw[<->, thick] (2.5,0) -- (m.west);
\end{tikzpicture}
\captionof{figure}{The imitation game (slide 22): the interrogator only sees written answers and
must guess which partner is the machine.}
\end{center}

\begin{nbintuition}[Behaviour instead of definitions]
Turing's trick is to judge intelligence by \emph{observable behaviour}, not by what happens
inside. The agent definition of AI (Section~\ref{sec:ai-definition}) has the same spirit.
The slides add a comic (``Turing test 2208'') where the human fails: ``the machine realised
immediately that it was me being the dumb human''. Today large language models make the
classic Turing test much less decisive than it looked in 1950.
\end{nbintuition}

\paragraph{An open parenthesis (slides 19--21).} What does it exactly mean to ``think'', to
``perceive'', to be ``conscious of being alive''? The slides link to the neuroscientist
Anil Seth (TED talk \emph{Your brain hallucinates your conscious reality}): according to him,
perception is a \emph{controlled hallucination}, the brain's best guess about the causes of
its sensory signals, constantly corrected by the senses. If we cannot define these words
precisely for humans, judging them in machines is even harder. (Closed parenthesis.)

\subsubsection{1956: the Dartmouth workshop and the founding fathers}
\slidesrc{Thinking machines: 24--33}

In summer 1956 a workshop at Dartmouth College (USA) coined the term
\textbf{``Artificial Intelligence''}. Its proposal contained the founding conjecture of the
field:

\begin{center}
\begin{minipage}{0.88\linewidth}
\itshape ``Every aspect of learning or any other feature of intelligence can in principle be
so precisely described that a machine can be made to simulate it.''
\end{minipage}
\end{center}

The slides list the participants (the ``founding fathers of AI'') and, step by step, what
each of them is remembered for:

\begin{center}
\small
\begin{tabularx}{\linewidth}{@{}l X@{}}
\toprule
\textbf{Person} & \textbf{Known for} \\
\midrule
John McCarthy & the \emph{Advice Taker} (a program reasoning with logical sentences, which
                inspired Prolog), the Lisp language, ALGOL \\
Marvin Minsky & artificial neural networks, the \emph{society of mind} (the mind as a
                society of simple interacting agents: an ancestor of multi-agent systems!) \\
Claude Shannon & information theory, cryptography, natural language processing \\
Ray Solomonoff & machine learning, prediction, Bayesian frameworks \\
Allen Newell \& Herbert Simon & the \emph{Logic Theory Machine} (proved theorems of
                \emph{Principia Mathematica}) and the \emph{General Problem Solver} \\
Arthur Samuel & search, alpha-beta pruning, machine learning (a checkers program that
                learned to play better than its author) \\
\midrule
Oliver Selfridge & Pandemonium theory (pattern recognition by many competing ``demons'') \\
Nathaniel Rochester & IBM, designer of the IBM 701 computer \\
Trenchard More & natural deduction (the proof system of Section~\ref{sec:natural-deduction-prop}) \\
\bottomrule
\end{tabularx}
\captionof{table}{The founding fathers of AI and their contributions (slides 26--33).}
\end{center}

\begin{nbintuition}[Symbolic and sub-symbolic, from day one]
Already at Dartmouth both souls of AI were present: logic and symbols (McCarthy, Newell and
Simon, More) and learning and neural networks (Minsky, Solomonoff, Samuel). The rest of the
history is a pendulum between the two.
\end{nbintuition}

\subsubsection{1966--1972: Shakey the robot}
\slidesrc{Thinking machines: 34--36}

\textbf{Shakey} was the first general-purpose mobile robot able to \emph{reason about its own
actions}. Other robots had to be instructed on each individual step; Shakey could analyse a
command and break it down into basic chunks by itself. The project combined robotics,
computer vision and natural language processing, and was the first to meld \emph{logical
reasoning} and \emph{physical action}. It was developed at the AI Center of the Stanford
Research Institute (now SRI International) from about 1966 to 1972 (project manager Charles
Rosen), programmed mainly in LISP, and its planning component was the \textbf{STRIPS}
planner. As the first logical, \textbf{goal-based agent}, Shakey lived in a limited world
(a few rooms with boxes).

\begin{nbexample}{what ``breaking a command into chunks'' means}
\notinslides
Command: ``push the box from room A to room B''. A planner like STRIPS knows actions with
preconditions and effects, e.g.\ \code{goto(Place)}, \code{push(Box, From, To)}, and finds
the sequence: \code{goto(room\_A)} $\to$ \code{goto(box)} $\to$
\code{push(box, room\_A, door)} $\to$ \code{push(box, door, room\_B)}. Your navigation app
and a warehouse robot do the same kind of planning today. Shakey is a deliberative agent
in the sense of Section~\ref{sec:agent-architectures}.
\end{nbexample}

\subsubsection{1970--1972: Prolog}
\slidesrc{Thinking machines: 37}

In 1970 Alain Colmerauer designed the \textbf{Prolog} language (from the French
\emph{programmation en logique}); the first Prolog system was implemented by Colmerauer and Philippe Roussel in 1972
in Marseille. Prolog is the language of Part~IV of these notes: a program is a set of
logical facts and rules, and running it means proving a query.

\subsubsection{The AI winters and the Lisp machines}
\slidesrc{Thinking machines: 38--43}

Early AI promised a lot. When the promises were not kept, funding collapsed: these periods
are called \textbf{AI winters}.

\begin{itemize}
  \item \textbf{First AI winter (1974--1980).} The \textbf{Lighthill report} (\emph{Artificial
        Intelligence: A General Survey}, Sir James Lighthill, 1973, for the British Science
        Research Council) gave a very pessimistic evaluation: ``in no part of the field have
        discoveries made so far produced the major impact that was then promised''. It formed
        the basis for the British government's decision to end support for AI research in
        all but three universities (Edinburgh, Sussex, Essex). Its main technical criticism:
        AI had failed to address the \textbf{combinatorial explosion} that appears when
        problems are taken from toy worlds to real-world domains.
  \item \textbf{AI renaissance (1980).} Symbolics commercialises \textbf{Lisp machines}:
        computers with hardware support for Lisp operations (data type testing, CDR coding,
        incremental garbage collection), running large Lisp programs very efficiently, with
        operating systems written entirely in Lisp. They powered the boom of
        \emph{expert systems} (rule-based programs encoding the knowledge of human experts).
  \item \textbf{Second AI winter (1987--1993).} With the microcomputer revolution, cheap
        desktop PCs soon ran Lisp programs even faster than Lisp machines, without special
        hardware. The Lisp machine market collapsed: most makers went out of business by the
        early 1990s, and AI lost credibility again.
\end{itemize}

\begin{nbexample}{combinatorial explosion in numbers}
\notinslides
Chess has about 35 legal moves per position. Looking $d$ moves ahead means examining about
$35^d$ positions:
\[
35^2 = 1\,225, \qquad 35^4 = 1\,500\,625, \qquad 35^6 \approx 1.8 \times 10^{9}, \qquad
35^{10} \approx 2.8 \times 10^{15}.
\]
Each extra level multiplies the work by 35. An algorithm that works on a toy problem can
become hopeless on a slightly bigger one: this is what the Lighthill report meant. Deep Blue
(1997) beat it on chess with special hardware plus pruning (alpha-beta, Samuel's legacy).
\end{nbexample}

\subsubsection{1997--2016: from chess to Go, the data-driven era}
\slidesrc{Thinking machines: 44--67}

\begin{itemize}[nosep]
  \item \textbf{1997}: IBM \textbf{Deep Blue} defeats the world chess champion Garry
        Kasparov (symbolic AI: search plus hand-written evaluation).
  \item \textbf{2008}: Google releases a \textbf{speech-to-search} iPhone app: you speak, the
        recording is sent to Google's servers, and the results come back in seconds.
  \item \textbf{2011}: IBM \textbf{Watson} defeats the champion Ken Jennings at the quiz show
        \emph{Jeopardy!} (natural language question answering).
  \item \textbf{2011/2012}: Apple introduces \textbf{Siri}, an intelligent personal assistant.
  \item \textbf{2012}: after seeing 10 million images from YouTube videos in three days, the
        16{,}000-computer network of \textbf{Google Brain} (one billion connections) began to
        recognise cats, even though it had never been taught what a cat looked like
        (unsupervised learning, 75\% accuracy).
  \item \textbf{2013}: China's Tianhe-2 doubles the world's top supercomputing speed (33.86
        petaflops); \textbf{2016}: Sunway TaihuLight triples it (93 petaflops). Raw computing
        power is one of the reasons for the success of deep learning.
  \item \textbf{2014}: Facebook's \textbf{DeepFace} identifies human faces in images: a
        nine-layer neural network with over 120 million connection weights, trained on four
        million images uploaded by Facebook users, said to be 97\% accurate (vs.\ 85\% for
        the FBI's Next Generation Identification system). It uses a 3D model to virtually
        rotate faces so that they look at the camera.
  \item \textbf{2015}: Google open-sources \textbf{TensorFlow}, a library for numerical
        computation with data flow graphs, now widely used for machine learning.
  \item \textbf{2016}: Google DeepMind's \textbf{AlphaGo} defeats Lee Sedol at Go, a game
        whose search space is far too large for Deep Blue-style brute force.
  \item \textbf{2016}: Uber pilots a \textbf{self-driving car} programme in Pittsburgh, PA.
\end{itemize}

\begin{nbwarning}[Two kinds of ``winning'']
Deep Blue and AlphaGo both beat a world champion, but in very different ways. Deep Blue
\emph{searched} millions of positions per second with rules written by humans (System~2
style, symbolic). AlphaGo \emph{learned} an intuition for good moves from millions of games
with neural networks, combined with a search (System~1 + System~2). Do not file both under
``AI wins at games'' without noting the difference: it is the theme of
Section~\ref{sec:symbolic-vs-subsymbolic}.
\end{nbwarning}

\subsection{Some more recent achievements}
\label{sec:recent-achievements}
\slidesrc{Thinking machines: 68--79}

After 2016 the pace accelerated. The slides show a selection of results; for each one,
here is what it is and why it matters.

\subsubsection{Games, medicine and emergent behaviour}
\slidesrc{Thinking machines: 68--71}

\begin{itemize}
  \item \textbf{AlphaStar} (DeepMind, 2019). StarCraft~II is a real-time strategy video game:
        many units, partial information (you do not see the whole map), real-time decisions.
        AlphaStar won its historic series against professional players (announced on
        January 24, 2019, after matches recorded in December 2018; e.g.\ 5--0 against the pro
        player MaNa).
  \item \textbf{Lung cancer screening} (\emph{Nature Medicine}, May 2019). A
        three-dimensional deep learning model analyses low-dose chest CT scans end to end and
        detects lung cancer with performance comparable to, or better than, expert
        radiologists.
  \item \textbf{Emergent tool use from multi-agent autocurricula} (OpenAI, ICLR 2020). Agents
        play \emph{hide-and-seek} in a simulated world with movable boxes and ramps. Nobody
        tells them how to use the objects; just by competing, the teams discover a sequence
        of strategies: hiders build shelters with boxes, seekers use ramps to jump in,
        hiders then lock the ramps away, seekers ``surf'' on boxes, and so on.
\end{itemize}

\begin{nbwhy}[A multi-agent system in action]
The hide-and-seek experiment is a \emph{multi-agent system}: several autonomous agents,
in a shared environment, competing (hiders vs.\ seekers) and cooperating (inside a team).
Complex behaviour \emph{emerges} from simple goals and interaction, with no central
controller. Keep this example in mind for Part~III (agents and MAS) and for agent-based
simulation (Section~\ref{sec:abms}).
\end{nbwhy}

\subsubsection{Self-driving cars: the SAE levels of driving automation}
\slidesrc{Thinking machines: 72--74}

How ``autonomous'' is an autonomous car? The SAE J3016 standard defines six levels.

\begin{center}
\small
\renewcommand{\arraystretch}{1.15}
\begin{tabularx}{\linewidth}{@{}c >{\raggedright\arraybackslash}X >{\raggedright\arraybackslash}X >{\raggedright\arraybackslash}X@{}}
\toprule
\textbf{Level} & \textbf{What the human must do} & \textbf{What the features do} & \textbf{Examples} \\
\midrule
\rowcolor{nbblue!6} 0 & You \textbf{are driving}, even if features are on; constantly supervise & warnings and momentary assistance & automatic emergency braking, blind-spot warning, lane departure warning \\
\rowcolor{nbblue!6} 1 & same as level 0 & steering \textbf{or} brake/acceleration support & lane centering \textbf{or} adaptive cruise control \\
\rowcolor{nbblue!6} 2 & same as level 0 & steering \textbf{and} brake/acceleration support & lane centering \textbf{and} adaptive cruise control together \\
\rowcolor{nbgreen!8} 3 & You are \textbf{not driving} while features are engaged, but you must take over when the system requests it & drives the vehicle under limited conditions & traffic jam chauffeur \\
\rowcolor{nbgreen!8} 4 & You are not driving; you will never be required to take over & drives under limited conditions, never needs you & local driverless taxi; pedals and steering wheel may not be installed \\
\rowcolor{nbgreen!8} 5 & You are not driving & drives everywhere, in all conditions & as level 4, but anywhere \\
\bottomrule
\end{tabularx}
\captionof{table}{SAE levels of driving automation (slide 72). Blue rows: \emph{driver support}
features (the human drives). Green rows: \emph{automated driving} features (the system
drives).}
\end{center}

The slides note (2024 news) that a \textbf{Level~3} car can now be bought in the US: the
Mercedes-Benz system works only on a small selection of roads, at limited speed and only in
some states, and its wide adoption will require regulators to catch up.

\begin{nbintuition}[Levels of autonomy, again later]
The question ``how much can the system decide by itself, and when must a human step in?''
is the question of \textbf{autonomy}, the defining property of agents
(Section~\ref{sec:agents-and-mas}). The jump from level 2 to level 3 is exactly the moment
in which responsibility moves from the human to the agent.
\end{nbintuition}

\subsubsection{Generative AI, protein folding and agentic AI}
\slidesrc{Thinking machines: 75--79}

On slide 75 the professor asked ChatGPT itself: ``which are the three most impressive
achievements in the AI research field in the last three years (short answer)''. The answer:
\begin{enumerate}[nosep]
  \item \textbf{Generative AI}: models such as GPT-4 and DALL-E generate coherent text,
        images and creative content.
  \item \textbf{AI in protein folding}: DeepMind's AlphaFold accurately predicts protein
        structures, solving a 50-year-old problem in biology.
  \item \textbf{AI alignment and safety research}: progress in interpretability, alignment
        and safety (e.g.\ by OpenAI and Anthropic) to control the behaviour of advanced AI.
\end{enumerate}

\begin{nbdef}{Protein folding problem}
A protein is a chain of amino acids. The \textbf{protein folding problem} is the question of
how a protein's amino acid sequence dictates its three-dimensional atomic structure. It
first emerged around 1960, with the first atomic-resolution protein structures.
\end{nbdef}

\textbf{AlphaFold2} (DeepMind) predicts the 3D structure of proteins from their amino acid
sequence with near-experimental accuracy; more than 200 million predicted structures have
been released. Knowing a protein's shape tells biologists how it works, which speeds up
drug discovery and protein design (2023 review; PNAS 2024, ``AlphaFold two years on'').

Finally, slide 79: \textbf{agentic AI}, presented as ``the next evolution of artificial
intelligence''. An agentic AI system does not just answer one question: given a goal, it
plans, calls tools (search engines, APIs, other programs), observes the results and acts
again, autonomously, over many steps.

\begin{nbwhy}[Old ideas, new engines]
Agentic AI is the 2020s name for an idea this course has studied for decades:
\textbf{autonomous agents} that perceive, decide and act towards goals, and
\textbf{multi-agent systems} in which several of them coordinate. The difference is the
engine: a large language model instead of (or together with) explicit beliefs, goals and
plans. Understanding classic agents (BDI, communication protocols, platforms) is what makes
it possible to engineer agentic AI systems that are reliable.
\end{nbwhy}

\subsection{But\ldots{} risks, and the answers to them}
\label{sec:risks-and-regulation}
\slidesrc{Thinking machines: 80--87}

The slides now switch tone (``But\ldots''): AI systems already take or support decisions
that affect people, and things can go wrong.

\subsubsection{Things that went wrong, or might}
\slidesrc{Thinking machines: 80--84}

\begin{center}
\small
\begin{tabularx}{\linewidth}{@{}l >{\raggedright\arraybackslash}X >{\raggedright\arraybackslash}X@{}}
\toprule
\textbf{Case} & \textbf{What happened} & \textbf{The underlying risk} \\
\midrule
State v.\ Loomis (2016--17) & The Wisconsin Supreme Court allowed judges to use an
  algorithmic risk-assessment score (predicting the risk of re-offending) in sentencing,
  but required a written warning about its limits. The defendant could not inspect how the
  proprietary algorithm worked. & opaque decisions about people's freedom; possible bias;
  no way to contest the reasoning \\
``Smart courts'' in China (2019) & Courts adopt AI assistants (e.g.\ Shanghai's
  ``System 206'') to check evidence and support judges. & automation of justice; who is
  accountable? \\
ChatGPT banned in Italy (2023) & The Italian data protection authority temporarily
  blocked ChatGPT over privacy concerns (personal data used for training, no age checks). &
  privacy and data protection \\
AI leaders at the UN Security Council & The heads of OpenAI (Sam Altman) and Anthropic
  (Dario Amodei) gave separate briefings on AI safety to the UN Security Council. &
  global safety and security \\
Incidents with AI agents (slide 84) & News reports about security incidents involving AI
  platforms and AI agents that were misused or hijacked. & autonomous agents acting in the
  real world can be attacked or misbehave \\
\bottomrule
\end{tabularx}
\captionof{table}{The cautionary examples of slides 80--84.}
\end{center}

\begin{nbintuition}[Why symbolic AI helps here]
In State v.\ Loomis the core problem is that nobody could \emph{see} why the system gave a
certain score. A symbolic system (a set of logical rules, a Prolog program, a BDI agent
with explicit plans) can \emph{explain} its conclusions: ``I concluded $X$ because of rule
$R$ applied to facts $F_1$ and $F_2$''. Explainability, verifiability and accountability
are among the strongest arguments for the symbolic side of AI.
\end{nbintuition}

\subsubsection{And hence: principles and guidelines}
\slidesrc{Thinking machines: 85--87}

\begin{itemize}
  \item \textbf{Asilomar AI Principles} (Future of Life Institute, 2017): 23 principles
        developed at the 2017 Asilomar conference, grouped into research issues, ethics and
        values (e.g.\ safety, transparency, responsibility, human control, privacy), and
        longer-term issues.
  \item \textbf{Ethics Guidelines for Trustworthy AI} (European Commission, High-Level
        Expert Group on AI, 8 April 2019). Trustworthy AI must be \emph{lawful},
        \emph{ethical} and \emph{robust}, and should meet seven key requirements (listed
        below, not in the slides).
  \item \textbf{Calls to slow down} (New York Times, September 2026): top AI leaders, among
        them the head of Anthropic, argued that capabilities are advancing so quickly that
        greater safety controls are needed across the industry.
\end{itemize}

\begin{nbfact}{The seven requirements for trustworthy AI (EU, 2019)}
\notinslides
\begin{enumerate}[nosep]
  \item human agency and oversight;
  \item technical robustness and safety;
  \item privacy and data governance;
  \item transparency (including explainability);
  \item diversity, non-discrimination and fairness;
  \item societal and environmental well-being;
  \item accountability.
\end{enumerate}
These guidelines were followed by the EU \emph{AI Act} (in force since 2024), which
regulates AI systems according to their level of risk.
\end{nbfact}

\subsection{Symbolic vs.\ sub-symbolic AI, and what this course is about}
\label{sec:symbolic-vs-subsymbolic}
\slidesrc{Thinking machines: 88--97}

The deck ends by going ``back to the first slides'': Kahneman's two systems
(Section~\ref{sec:system1-system2}) give the key to read the whole history of AI.

\subsubsection{Two families of AI}
\slidesrc{Thinking machines: 89--90}

\begin{center}
\begin{tikzpicture}[font=\scriptsize, >=Stealth,
  c/.style={circle, fill=#1, text=white, minimum size=1.55cm, align=center, inner sep=1pt}]
  % sub-symbolic
  \node[font=\small\bfseries] at (3.1,2.75) {Sub-symbolic AI};
  \node[c=nborange] (n1) at (0,0) {Numerical\\represen-\\tation};
  \node[c=nborange] (n2) at (2.05,0) {Machine\\learning};
  \node[c=nborange] (n3) at (4.1,0) {Distributed\\represen-\\tation};
  \node[font=\large\bfseries, inner sep=0pt] at (1.025,0) {+};
  \node[font=\large\bfseries, inner sep=0pt] at (3.075,0) {+};
  \node[font=\large\bfseries, inner sep=0pt] at (5.125,0) {=};
  \node[c=nborange!70!black] at (6.15,0) {The sub-\\symbolic\\approach};
  % symbolic
  \begin{scope}[xshift=10.1cm]
  \node[font=\small\bfseries] at (0,2.75) {Symbolic AI (GOFAI)};
  \node[c=nbblue!70!black] (s) at (0,-0.6) {The\\symbolic\\approach};
  \node[c=nbblue] (k) at (0,1.0-1.6+2.2) {Knowledge\\represen-\\tation};
  \node[c=nbblue] (p) at (-1.55,-1.95) {Problem-\\solving\\algorithms};
  \node[c=nbblue] (l) at (1.55,-1.95) {Logical\\reasoning};
  \draw[->, very thick, nbgray] (s) -- (k);
  \draw[->, very thick, nbgray] (s) -- (p);
  \draw[->, very thick, nbgray] (s) -- (l);
  \end{scope}
\end{tikzpicture}
\captionof{figure}{The two families of AI (slides 89--90). Sub-symbolic AI: knowledge is
numbers spread over many parameters, learned from data. Symbolic AI: knowledge is explicit
symbols, manipulated by reasoning and search.}
\end{center}

\begin{nbdef}{Sub-symbolic AI}
AI where knowledge is encoded as \textbf{numbers} (e.g.\ the weights of a neural network),
\textbf{learned} from data (machine learning), and \textbf{distributed}: no single number
means ``cat'' or ``stop sign''; the meaning is spread over millions of parameters.
\end{nbdef}

\begin{nbdef}{Symbolic AI (``Good Old-Fashioned AI'', GOFAI)}
AI where knowledge is written explicitly as \textbf{symbols} that humans can read (facts,
rules, logical formulas, plans), and intelligence is obtained by \textbf{knowledge
representation}, \textbf{logical reasoning} and \textbf{problem-solving algorithms}
(search, planning).
\end{nbdef}

\begin{nbexample}{the same knowledge, two ways}
Knowledge: ``Socrates is a human; all humans are mortal''.
\begin{itemize}[nosep]
  \item \emph{Symbolic}: two lines of Prolog (Part~IV),
        \code{human(socrates).} and \code{mortal(X) :- human(X).}
        Asking \code{?- mortal(socrates).} gives \code{true}, and the system can show the
        proof: rule + fact.
  \item \emph{Sub-symbolic}: a language model has read millions of sentences about
        Socrates and mortality; asked ``is Socrates mortal?'' it answers ``yes'' because
        this is the most likely continuation. The ``knowledge'' is somewhere in billions of
        weights and cannot be pointed at.
\end{itemize}
\end{nbexample}

\subsubsection{Thinking fast and slow in AI}
\slidesrc{Thinking machines: 88, 91}

The paper \emph{Thinking Fast and Slow in AI} (G.~Booch, F.~Rossi and others, AAAI 2021)
draws a loose parallel between Kahneman's two systems and the two main lines of work in AI:

\begin{center}
\small
\begin{tabularx}{\linewidth}{@{}l >{\raggedright\arraybackslash}X >{\raggedright\arraybackslash}X@{}}
\toprule
 & \textbf{System 1 $\approx$ machine learning} & \textbf{System 2 $\approx$ symbolic AI} \\
\midrule
Style & data-driven, fast, intuitive & knowledge-driven, slow, deliberate \\
Good at & perception from sensory data (seeing, reading, hearing); building (possibly
  imprecise and biased) models from data & complex problems: logic, search, optimisation,
  planning, with explicit knowledge, symbols and high-level concepts \\
Weak at & common-sense reasoning, basic notions of causality, explanations & noisy
  perception, learning from raw data, scaling to huge open-ended domains \\
Example & recognise a pedestrian in a camera image & decide that the car must yield
  because traffic rules say so \\
\bottomrule
\end{tabularx}
\captionof{table}{The parallel of slide 91 (``Thinking Fast and Slow in AI''), with examples.}
\end{center}

As in humans, the two sets of capabilities are \textbf{symbiotic}: this supports
\textbf{hybrid, neurosymbolic} approaches to AI.

\begin{nbdef}{Neurosymbolic AI}
\textbf{Neurosymbolic AI} combines neural networks (sub-symbolic learning) with symbolic
representations and reasoning, to obtain systems that are more versatile and more
dependable than either alone (slides 92--95: articles in \emph{Communications of the ACM}
2022, IBM Research, the book \emph{Neuro-Symbolic Artificial Intelligence: The State of the
Art}, and the journal \emph{Neurosymbolic Artificial Intelligence}, founded in 2023).
\end{nbdef}

\begin{nbexample}{a neurosymbolic self-driving car}
\notinslides
\begin{center}
\begin{tikzpicture}[>=Stealth, font=\small, node distance=0.6cm,
  b/.style={draw, rounded corners, thick, align=center, minimum height=0.9cm, text width=2.7cm}]
  \node[b, draw=nbgray] (cam) {camera image};
  \node[b, draw=nborange, fill=nborange!8, right=of cam] (nn) {\textbf{neural network}\\(System 1)};
  \node[b, draw=nbblue, fill=nbblue!8, below=0.9cm of nn] (rules) {\textbf{traffic rules}\\(System 2, symbolic)};
  \node[b, draw=nbgreen, right=of rules] (act) {action: brake};
  \draw[->, thick] (cam) -- (nn);
  \draw[->, thick] (nn) -- node[right, font=\scriptsize, align=left]{facts: \code{pedestrian(p1)},\\\code{on\_crossing(p1)}} (rules);
  \draw[->, thick] (rules) -- (act);
\end{tikzpicture}
\end{center}
The neural part turns pixels into \emph{symbols} (``there is a pedestrian on the zebra
crossing''); the symbolic part applies explicit, auditable rules
(``if a pedestrian is on a crossing, then yield''). If the car brakes, we can explain why.
\end{nbexample}

\subsubsection{What SDAI is about}
\slidesrc{Thinking machines: 96--97}

\begin{nbfact}{The course in one paragraph (slide 96)}
The Symbolic and Distributed AI course is about the \textbf{symbolic component of
neurosymbolic AI}. It exploits symbolic AI in the context of \textbf{distributed AI
systems}, where (symbolic, but not only) intelligent software \textbf{agents} distributed
across different machines operate \textbf{autonomously}, to \textbf{coordinate/collaborate}
in order to achieve some common goal, or to \textbf{compete} to access shared resources. A
system consisting of these distributed agents, as well as their environment and the
organisations they are part of, is named \textbf{multi-agent system} (MAS).
\end{nbfact}

\begin{center}
\begin{tikzpicture}[>=Stealth, font=\small,
  b/.style={draw, rounded corners, thick, align=center, text width=3.4cm, minimum height=1cm}]
  \node[b, draw=nbblue, fill=nbblue!8] (sym) at (0,0) {\textbf{Symbolic AI}\\logic (Part II)\\Prolog (Part IV)};
  \node[b, draw=nbteal, fill=nbteal!8] (dis) at (7,0) {\textbf{Distributed AI}\\agents, MAS (Part III)};
  \node[b, draw=nbpurple, fill=nbpurple!8, text width=5.6cm] (mas) at (3.5,-2.2)
    {\textbf{Symbolic agents in a MAS}\\BDI agents, AgentSpeak/Jason, communication,
     JADE, simulation (Part V)};
  \draw[->, thick] (sym) -- (mas);
  \draw[->, thick] (dis) -- (mas);
\end{tikzpicture}
\captionof{figure}{How the parts of these notes fit together.}
\end{center}

Welcome on board, and let's have fun with S\&DAI!


\begin{nbsummary}[Introduction to AI]
\begin{itemize}[nosep]
  \item \textbf{Intelligence}: general capability to reason, plan, solve problems, think
        abstractly, comprehend complex ideas, learn quickly and learn from experience. Many
        kinds exist (Gardner's 9 intelligences, emotional intelligence).
  \item \textbf{System 1} = fast, automatic, intuitive; \textbf{System 2} = slow, effortful,
        logical (Kahneman, 2011). Parallel in AI: System~1 $\approx$ machine learning,
        System~2 $\approx$ symbolic AI (logic, planning, search).
  \item \textbf{AI} (Russell \& Norvig): the study of \emph{intelligent agents}, devices
        that perceive their environment and act to maximise the chance of achieving their goals.
  \item \textbf{Key dates}: 1950 Turing test; 1956 Dartmouth workshop (the term ``AI'');
        1966--72 Shakey; 1970--72 Prolog; 1974--80 and 1987--93 AI winters; 1997 Deep Blue;
        2011 Watson; 2012 Google Brain cats; 2016 AlphaGo; 2019 AlphaStar; 2020+ AlphaFold,
        generative AI; now agentic AI.
  \item \textbf{Risks}: biased decisions (State v.\ Loomis), privacy (ChatGPT blocked in
        Italy), safety and security of autonomous agents. \textbf{Responses}: Asilomar
        principles (2017), EU Ethics guidelines for trustworthy AI (2019).
  \item \textbf{Sub-symbolic AI} = numerical representation + machine learning + distributed
        representation. \textbf{Symbolic AI} (GOFAI) = knowledge representation + logical
        reasoning + problem-solving algorithms. \textbf{Neurosymbolic AI} combines them.
  \item \textbf{SDAI} = the symbolic component of neurosymbolic AI, used inside distributed
        systems of autonomous agents: a \textbf{multi-agent system (MAS)}.
\end{itemize}
\end{nbsummary}
